La struttura algebrica $(\mathbb{Z}, +)$ risulta quindi un gruppo abeliano.
Una generica struttura algebrica $(\mathbb{G}, \cdot)$, formata da un insieme
$\mathbb{G} \neq \emptyset$ e un'operazione binaria "$\cdot$", è un gruppo se e solo se
soddisfa le seguenti proprietà.

\begin{property}
    \textbf{Chiusura di $\mathbb{G}$ su $\cdot$}
    \[\forall a, b \in \mathbb{G}, \; (a \cdot b) \in \mathbb{G}.\]
\end{property}
\begin{property}
    \textbf{Associativa}
    \[\forall a, b, c \in \mathbb{G}, \; (a \cdot b) \cdot c = a \cdot (b \cdot c).\]
\end{property}
\begin{property}
    \textbf{Esistenza dell'elemento neutro}
    \[\forall a \in \mathbb{G}, \; \exists e \in \mathbb{G} \mid a \cdot e = e \cdot a = a.\]
\end{property}
\begin{property}
    \textbf{Esistenza dell'elemento inverso}
    \[\forall a \in \mathbb{G}, \; \exists b \in \mathbb{G} \mid a \cdot b = b \cdot a = e.\]
\end{property}

\begin{observation}
    Non è necessaria la commutatività dell'operazione binaria per definire un gruppo
\end{observation}

Un gruppo si dice inoltre commutativo o abeliano se e solo se soddisfa la seguente proprietà.
\begin{property}
    \textbf{Commutativa}
    \[\forall a, b \in \mathbb{G}, \; a \cdot b = b \cdot a.\]
\end{property}

Affinché una generica struttura algebrica sia un gruppo abeliano,
è dunque necessario che soddisfi contemporaneamente tutte le proprietà sopra elenecate.